

\section{Ranking from pairwise comparisons}
\label{sec:ranking}




The ranking task---which seeks to identify a consistent ordering of several items based on
(partially) revealed preference information about them---is encountered in numerous contexts
including web search, crowd sourcing, social choice, peer grading, and so on \citep{dwork2001rank,chen2013pairwise,caplin1991aggregation,shah2013case}.
Of particular interest is the ``preference-based'' observation model,
in which we are only given relative comparisons of a few items (as opposed to individual scores of them).
In practice, comparison data of this kind abound,
partly because humans often find it easier to make a preference over two or a couple of items than to assign specific ratings to many individual ones.
The emergence of crowdsourcing platforms such as Amazon Mechanical Turk further widens the availability of comparison data,
where binary judgments over pairs of items are often solicited from a pool of non-experts.
In this section, we concentrate on pairwise comparisons and
explore the potential of spectral methods for the ranking task.




\subsection{The Bradley-Terry-Luce model and assumptions \label{subsec:Problem-setup-and-ranking}}

To formulate the problem in a statistically sound manner, we introduce a classical parametric model,
called the Bradley-Terry-Luce (BTL) model~\citep{bradley1952rank,ford1957solution,luce2012individual},
to describe the generating process of pairwise comparisons.


\paragraph{Latent preference scores.}
Imagine that there are $n$ items to be ranked.
A key component of the BTL model is the assignment of a latent preference score to each item; more concretely, the BTL model hypothesizes on the existence of an unseen preference score vector
%
\begin{equation}
	\bm{w}^{\star}=[w_{1}^{\star},w_{2}^{\star},\cdots, w_{n}^{\star}]^{\top},
	\label{eq:ranking-score}
\end{equation}
%
with $w_{i}^{\star}>0$ assigned to the $i$-th item ($1\leq i\leq n$).
The ranks of these items are therefore determined exclusively by their (relative) preference scores: an item with a larger score is ranked higher.
Throughout this section, we denote by $\kappa$ a sort of  condition number as follows
%
\begin{equation}
	\kappa \coloneqq\frac{\max_{1\leq i\leq n}w_{i}^{\star}}{\min_{1\leq i\leq n}w_{i}^{\star}}.
	\label{eq:ranking-well-condition}
\end{equation}

%$[w_{\min}, w_{\max}]$ the dynamic range of the scores (so that  $w_i[w_{\min}, w_{\max}]$ for all $1\leq i\leq n$), and let



\paragraph{Pairwise comparisons.}
Equipped with the aforementioned score
vector, the BTL model posits that: the probability of an item winning a paired comparison is determined entirely by the
relative scores of the two items involved. To be precise, when comparing  every pair $(i,j)$ of items, the model assumes that
%
\begin{equation}
	\mathbb{P} \big\{ \text{item }j\text{ is preferred over item }i \big\}
	= \frac{w_{j}^{\star}}{w_{i}^{\star}+ w_{j}^{\star}},
	\label{eq:BTL}
\end{equation}
%
asserting that an item assigned a higher preference score is more likely to win.
In this section, we assume access to a comparison between every pair of items. To be precise, for each pair
$(i,j)$ ($1\le i<j\le n$), we observe an independent binary comparison outcome $y_{i,j}$ following the BTL model \eqref{eq:BTL}:
%
\[
y_{i,j}=\begin{cases}
1, & \text{with probability }\frac{w_{j}^{\star}}{w_{i}^{\star}+w_{j}^{\star}},\\
0, & \text{otherwise},
\end{cases}
\]
%
where $y_{i,j}=1$ means item $j$ beats item $i$ and $y_{i,j}=0$ otherwise.
By convention, we set $y_{i,j}=1-y_{j,i}$ for all $i>j$.
%
%To simplify the notation hereafter, we denote
%%
%\begin{equation}
%	y_{i,j}^{\star}\coloneqq\frac{w_{j}^{\star}}{w_{i}^{\star}+w_{j}^{\star}}.
%	\label{eq:true-comparison-prob}
%\end{equation}
%%
%

%

\paragraph{Goal. }
With the BTL parametric model in mind,
a natural strategy is to start by estimating the underlying scores $\{ w_i^{\star}\}$ based on the pairwise comparisons in hand, followed by a ranking step performed in  accordance with the estimated scores.
In this section, we shall focus on characterizing the statistical accuracy of spectral methods in accomplishing the meta task of preference score estimation,
and will remark in passing on the ranking step that follows. Obviously, from \eqref{eq:BTL}, we can only hope for estimating $\{ w_i^{\star}\}$ up to some global scaling ambiguity.


% Motivated by this idea, this section explores the capability of spectral methods in accomplishing the meta task of score estimation.
%
%Due to the scale invariance of the BTL observation model (see~(\ref{eq:BTL})),
%it is assumed without loss of generality that
%%
%\begin{equation}
% 	\sum_{i=1}^{n}w_{i}^{\star}=1.
%	\label{eq:BTL-normalization}
%\end{equation}
%%



\subsection{A spectral ranking algorithm \label{subsec:Spectral-method-for-ranking}}

At first glance, the recipe we have introduced for designing  spectral methods seems to have no direct bearing on the BTL model.
Somewhat unexpectedly, a closer inspection unveils an intimate connection between the BTL model and a reversible Markov chain,
whose stationary distribution embodies crucial information about the score vector of interest.
This in turn lays a solid foundation for the spectral algorithm described below, originally developed by \citet{negahban2016rank}.

The first step is to convert the pairwise comparison data
$\{y_{i,j}\}_{i\neq j}$ into a probability transition matrix $\bm{P}=[P_{i,j}]_{1\leq i,j\leq n}$, in a way that
%
\begin{equation}
P_{i,j}=\begin{cases}
\frac{1}{n}y_{i,j}, & \text{if }i\neq j,\\
1-\sum_{j:j\neq i}\frac{1}{n}y_{i,j}, & \text{otherwise}.
\end{cases}
\label{defn:ranking-P-empirical}
\end{equation}
%
By construction of $\bm{P}$, all of its entries are non-negative and the entries in each row sum up to one,
thus confirming that  $\bm{P}$ is a probability transition matrix.
The spectral algorithm then computes the leading left eigenvector $\bm{\pi}$ of  $\bm{P}$, returning it as
 the estimate for the underlying score vector $\bm{w}^{\star}$.

To make sense of the rationale behind this algorithm, it is helpful to
look at the mean $\bm{P}^{\star} =[P_{i,j}^{\star}]_{1\leq i,j\leq n} \coloneqq\mathbb{E}[\bm{P}]$, which obeys
%
\begin{equation}
P_{i,j}^{\star}=\begin{cases}
	\frac{1}{n} \frac{w_j^{\star}}{w_i^{\star}+w_j^{\star}}, & \text{if }i\neq j,\\
	1- \frac{1}{n}  \sum_{j:j\neq i}\frac{w_j^{\star}}{w_i^{\star}+w_j^{\star}}, & \text{otherwise}.
\end{cases}\label{defn:ranking-P-star}
\end{equation}
%
Clearly, this matrix $\bm{P}^{\star}$ is a probability transition matrix as well.
As can be straightforwardly verified, the vector $\bm{\pi}^{\star}=[\pi_{i}^{\star}]_{1\leq i\leq n}$ defined by
%
\begin{align}
	%\pi_i^{\star} = \frac{w_i^{\star}}{\sum_{l=1}^n w_l^{\star}},\quad 1\leq i\leq n
	%\qquad \text{or} \qquad
	\bm{\pi}^{\star} = \frac{1}{\bm{1}^{\top}\bm{w}^{\star}} \bm{w}^{\star}
	\label{defn:pi-i-star-BTL}
\end{align}
%
satisfies the following conditions:
%
\begin{itemize}
	\item $\bm{\pi}^{\star}$ is a probability vector (i.e., $\pi_i^{\star}\geq 0$ for all $i$ and $\sum_i\pi_i^{\star}=1$);
	\item $\bm{\pi}^{\star}$ satisfies the detailed balance equations as follows:
%
\[
	\pi_{i}^{\star}P_{i,j}^{\star} = \pi_{j}^{\star}P_{j,i}^{\star},\qquad\text{for all }(i,j).
\]
%
\end{itemize}
%
Classical Markov chain theory \citep{bremaud2013markov} thus tells us that
$\bm{P}^{\star}$ represents a reversible Markov chain, whose stationary distribution is precisely given by $\bm{\pi}^{\star}$ (this can easily be verified using the definition of the stationary distribution) and corresponds to the normalized preference scores.
As a consequence, we hold the intuition that:  $\bm{\pi}$ is close to $\bm{\pi}^{\star}$---and hence $\bm{w}^{\star}$ up to some global scaling---as long as $\bm{P}$ approximates $\bm{P}^{\star}$ reasonably well.






\subsection{Performance guarantees}
\label{sec:theory-ranking}


This subsection develops theoretical support for the above spectral ranking algorithm,
based on the eigenvector perturbation theory developed previously for probability transition matrices in Section~\ref{sec:eigenvector-theory-DK}.
For notational convenience, we shall use $\bm{E}\coloneqq\bm{P}-\bm{P}^{\star}$ to
denote the difference of the above two transition matrices of interest.


By virtue of Theorem~\ref{thm:DK_asym},
the perturbation of the stationary distribution of a reversible Markov
chain $\bm{P}^{\star}$ is dictated by two important quantities: (i)
the spectral gap $1-\max\left\{ \lambda_{2}(\bm{P}^{\star}),-\lambda_{n}(\bm{P}^{\star})\right\} $,
and (ii) the noise size $\|\bm{E}\|_{\bm{\pi}^{\star}}$ (recall the definition of $\|\cdot\|_{\bm{\pi}^{\star}}$  in Section~\ref{subsec:Setup-and-notation-prob-matrix}).
These two quantities are controlled respectively via the following two lemmas, whose
proofs can be found in Section~\ref{subsec:Proof-of-auxilliary-ranking}.

\begin{lemma}
	\label{lemma:ranking-gap}
	Consider the settings and notation in Sections~\ref{subsec:Problem-setup-and-ranking} and~\ref{subsec:Spectral-method-for-ranking}. It follows that
%
\[
1-\max\big\{ \lambda_{2}(\bm{P}^{\star}),-\lambda_{n}(\bm{P}^{\star}) \big\} \geq\frac{1}{2\kappa^{2}},
\]
%
where we recall the definition of $\kappa$  in (\ref{eq:ranking-well-condition}).
\end{lemma}
%
\begin{lemma}
	\label{lemma:ranking-noise}
Consider the settings and notation in
Sections~\ref{subsec:Problem-setup-and-ranking} and \ref{subsec:Spectral-method-for-ranking}, and recall that $\bm{E}\coloneqq\bm{P}-\bm{P}^{\star}$.
With probability at least $1-O(n^{-8})$,
%
\[
	\|\bm{E}\|_{\bm{\pi}^{\star}}\leq \sqrt{\kappa}\,\|\bm{E}\| \lesssim\sqrt{\frac{\kappa\log n}{n}}.
\]
%
\end{lemma}
%
Now we are well prepared to assess the quality of the spectral estimate, as summarized below, whose proof is given at the end of this subsection.
%
\begin{theorem}
\label{thm:ranking}
Consider the settings and algorithm in Sections~\ref{subsec:Problem-setup-and-ranking}
and \ref{subsec:Spectral-method-for-ranking}. Suppose that $n\geq C\kappa^{5}\log n$
for some sufficiently large constant $C>0$. Then with probability
exceeding $1-O(n^{-8})$, one has
%
\begin{align}
	\frac{\|\bm{\pi}-\bm{\pi}^{\star}\|_{2}}{\|\bm{\pi}^{\star}\|_{2}}\lesssim\kappa^{2.5}\sqrt{\frac{\log n}{n}}.
	\label{eq:pi-pistar-quality-ranking}
\end{align}
%
\end{theorem}
%
Given the construction \eqref{defn:pi-i-star-BTL} of $\bm{\pi}^{\star}$, this theorem implies the existence of a scalar $z>0$ such that
%
\[
	\frac{\| z\bm{\pi} - \bm{w}^{\star}\|_{2}}{\|\bm{w}^{\star}\|_{2}}\lesssim\kappa^{2.5}\sqrt{\frac{\log n}{n}}
\]
%
holds with high probability. It is worth noting that one cannot possibly retrieve the global scaling factor $z$, due to the invariance of the BTL observation model under global scaling (cf.~(\ref{eq:BTL})).



To interpret the effectiveness of this theorem, consider, for example, the case when $\kappa=O(1)$ (so that all the latent scores
$w_{i}^{\star}$ are about the same order). Theorem~\ref{thm:ranking}
tells us that the relative estimation error of $\bm{\pi}$ is vanishing
as the number $n$ of items increases. As it turns out, this statistical error rate \eqref{eq:pi-pistar-quality-ranking}
 is near minimax-optimal up to a logarithmic
factor; see~\citet[Theorem 3]{negahban2016rank}. In fact, with a more careful
analysis, one can further eliminate this extra $\log n$ factor and establish (orderwise) minimax optimality of this algorithm,
as has been done in \citet[Theorem 5.2]{chen2017spectral}.

Caution needs to be exercised, however, that high score estimation accuracy alone does not necessarily imply appealing ranking accuracy.
For instance, if the goal is to identify the top-$K$ ranked items---a problem commonly referred to as ``top-$K$ ranking'' \citep{chen2015spectral}---then the ranking accuracy also relies heavily on the separation between the score of the $K$-th ranked item and that of the $(K+1)$-th ranked item (namely, whether the set of top-$K$ ranked items is sufficiently distinguishable from the remaining ones). Fortunately, the spectral ranking algorithm introduced in this section remains minimax optimal when it comes to top-$K$ ranking, through a refined $\ell_\infty$ perturbation theory to be introduced in Chapter~\ref{cha:Linf-theory}. The interested reader is referred to \citet{chen2017spectral} for details.


\begin{proof}[Proof of Theorem~\ref{thm:ranking}]
%Recall that $\bm{E}=\bm{P}-\bm{P}^{\star}$ denotes
%the noise matrix.
Invoke Theorem~\ref{thm:DK_asym} to see that
%
\begin{align}
	\|\bm{\pi}-\bm{\pi}^{\star}\|_{\bm{\pi}^{\star}}
	& \leq \frac{\big\Vert \bm{\pi}^{\star\top}\bm{E}\big\Vert _{\bm{\pi}^{\star}}}
	{1-\max\left\{ \lambda_{2}(\bm{P}^{\star}),-\lambda_{n}(\bm{P}^{\star})\right\} -\left\Vert \bm{E}\right\Vert _{\bm{\pi}^{\star}}} \nonumber\\
	& \leq 4 \kappa^{2}\big\Vert \bm{\pi}^{\star\top}\bm{E}\big\Vert _{\bm{\pi}^{\star}},
	\label{eq:pi-perturbation-UB1-ranking}
\end{align}
%
provided that
%
\begin{align}
	1-\max\big\{ \lambda_{2}(\bm{P}^{\star}),-\lambda_{n}(\bm{P}^{\star})\big\}
	-\left\Vert \bm{E}\right\Vert _{\bm{\pi}^{\star}}\geq1/(4\kappa^{2}) .
	\label{eq:condition-E-eigen-gap-bound-ranking}
\end{align}
%
From Lemma~\ref{lemma:ranking-gap} and Lemma~\ref{lemma:ranking-noise},
we know that Condition~\eqref{eq:condition-E-eigen-gap-bound-ranking} holds true with probability at least $1-O(n^{-8})$,
with the proviso that $n\geq C\kappa^{5}\log n$ for some sufficiently large constant $C>0$.


Additionally, letting $\pi_{\min}^{\star} \coloneqq \min_i\pi_i^{\star} $ and $\pi_{\max}^{\star} \coloneqq \max_i\pi_i^{\star} $,
we can easily see from the definition of $\|\cdot\|_{\bm{\pi}^{\star}}$ (i.e., $\|\bm{v}\|_{\bm{\pi}^{\star}}=\sqrt{\sum_i \pi_i^{\star}v_i^2 }$ for any vector $\bm{v}$)  that
%
\[
	\|\bm{v}\|_{\bm{\pi}^{\star}} \overset{\mathrm{(i)}}{\leq} \sqrt{\pi^{\star}_{\max }} \, \|\bm{v}\|_2 , \qquad \text{and} \qquad
	\|\bm{v}\|_2  \overset{\mathrm{(ii)}}{\leq}  \frac{1}{\sqrt{  \pi^{\star}_{\min} } } \, \|\bm{v}\|_{\bm{\pi}^{\star}} ,
\]
%
which allows us to further obtain
%
\begin{align*}
\|\bm{\pi}-\bm{\pi}^{\star}\|_{2} & \leq\frac{1}{\sqrt{\pi_{\min}^{\star}}}\|\bm{\pi}-\bm{\pi}^{\star}\|_{\bm{\pi}^{\star}}\leq\frac{4\kappa^{2}}{\sqrt{\pi_{\min}^{\star}}}\|\bm{\pi}^{\star\top}\bm{E}\|_{\bm{\pi}^{\star}}\leq4\kappa^{2.5}\|\bm{\pi}^{\star\top}\bm{E}\|_{2}\\
 & \leq4\kappa^{2.5}\|\bm{E}\|\,\|\bm{\pi}^{\star}\|_{2}.
\end{align*}
%
Here, the first inequality comes from (ii), the second inequality is a consequence of \eqref{eq:pi-perturbation-UB1-ranking}, whereas the third one results from (i).
The proof is then completed by applying the high-probability bound $\|\bm{E}\|\lesssim\sqrt{(\log n)/n}$ derived in Lemma~\ref{lemma:ranking-noise}.
\end{proof}





\subsection{Proof of auxiliary lemmas \label{subsec:Proof-of-auxilliary-ranking}}

Before delving into the proof, we state a general comparison theorem, which is attributed to \citet{diaconis1993comparison},
that relates the spectral gap of a reversible Markov chain with that
of another (possibly more tractable) reversible chain. We refer the interested reader
to \citet[Lemma 6]{negahban2016rank} for a proof of the following result.
%
\begin{lemma}
	\label{lemma:gap-comparison}
	Consider two reversible Markov chains over the state space $\{1,2,\cdots,n\}$.
	Let $\bm{P}$ and $\bm{\pi}$ (resp.~$\widetilde{\bm{P}}$ and $\widetilde{\bm{\pi}}$) denote
	the transition matrix and the stationary distribution of the first (resp.~second) chain.
	In addition, set
	%
\[
	\alpha\coloneqq\min_{i,j}
	\frac{\pi_i P_{i,j}}{\widetilde{\pi}_i \widetilde{P}_{i,j}},
	\qquad\text{and}\qquad
	\beta\coloneqq\max_i\frac{\pi_i}{\widetilde{\pi}_i}.
\]
%
Then one has
%
\[
	\frac{1-\max\left\{ \lambda_{2}(\bm{P}),-\lambda_{n}(\bm{P})\right\} }{1-\max\big\{ \lambda_{2}(\widetilde{\bm{P}}),-\lambda_{n}(\widetilde{\bm{P}})\big\} }
	\geq\frac{\alpha}{\beta}.
\]
%
\end{lemma}
%
Armed with this comparison lemma, we are ready to present the proof of Lemma \ref{lemma:ranking-gap}.




\paragraph{Proof of Lemma \ref{lemma:ranking-gap}.}
In order to control the spectral gap with the aid of Lemma~\ref{lemma:gap-comparison}, we construct an auxiliary transition matrix
%
\[
\bm{Q}^{\star}=\frac{1}{n}\bm{1}\bm{1}^{\top},
\]
%
which clearly corresponds to a reversible Markov chain with  stationary distribution $\bm{u}^{\star}=(1/n)\cdot\bm{1}$.  The eigengap of this newly constructed reversible Markov chain is
$$1-\max\left\{ \lambda_{2}(\bm{Q}^{\star}),-\lambda_{n}(\bm{Q}^{\star})\right\} =1,$$
since $\lambda_{2}(\bm{Q}^{\star})=\lambda_{n}(\bm{Q}^{\star})=0$.  Therefore, we only need to bound $\alpha$ and $\beta$.


Recalling the construction of $\bm{P}^{\star}$ in (\ref{defn:ranking-P-star}),
we can straightforwardly check that
%
\[
	\pi_{i}^{\star}P^{\star}_{i,j} %= \pi_{i}^{\star} \frac{1}{n}\frac{w_{j}^{\star}}{w_{i}^{\star}+w_{j}^{\star}}
	= \frac{1}{n}\frac{\pi_{i}^{\star}\pi_{j}^{\star}}{\pi_{i}^{\star}+\pi_{j}^{\star}} \geq\frac{1}{2n}\min\{\pi_{i}^{\star},\pi_{j}^{\star}\}
\]
%
for every $i\neq j$,  and in addition,
%
\[
\pi_{i}^{\star}P_{i,i}^{\star}=\pi_{i}^{\star}\Bigg[1-\sum_{j:j\neq i}\frac{1}{n}\frac{w_{j}^{\star}}{w_{i}^{\star}+w_{j}^{\star}}\Bigg]\geq\pi_{i}^{\star}\Bigg[1-\sum_{j:j\neq i}\frac{1}{n}\Bigg]=\frac{1}{n}\pi_{i}^{\star}.
\]
%
Combining the previous two inequalities, we obtain
%
\[
	\min_{i,j} \big( \pi^{\star}_i P^{\star}_{i,j} \big)
	\geq\frac{1}{2n}\min_{1\leq i\leq n}\pi_{i}^{\star}
	= \frac{1}{2n\kappa}\max_{1\leq i\leq n}\pi_{i}^{\star}
	\geq\frac{1}{2n^{2}\kappa},
\]
%
where the last relation holds since $\max_{ i}\pi_{i}^{\star}\geq \frac{1}{n}\sum_i\pi_i^{\star} = \frac{1}{n}$.
This together with the construction of $\bm{Q}^{\star}$ further leads to
%
\[
	\alpha \coloneqq \min_{i,j}\frac{\pi_{i}^{\star}P^{\star}_{i,j}}{u_{i}^{\star}Q^{\star}_{i,j}}
	=n^{2}\min_{i,j} \big( \pi^{\star}_i P^{\star}_{i,j} \big)
	\geq\frac{1}{2\kappa}.
\]
%
In regard to $\beta$, it is seen that
%
\[
	\beta\coloneqq \max_i \frac{\pi^{\star}_i}{u^{\star}_i} = n\max_{ i}\pi_{i}^{\star}
	= n\kappa \min_{i}\pi_{i}^{\star}
	\leq \kappa,
\]
%
where the final inequality follows since $\min_i\pi_i\leq \frac{1}{n}\sum_i\pi_i = \frac{1}{n}$. With the preceding bounds on $\alpha$ and $\beta$ in place,
Lemma~\ref{lemma:gap-comparison} informs us that
%
\[
\frac{1-\max\left\{ \lambda_{2}(\bm{P}^{\star}),-\lambda_{n}(\bm{P}^{\star})\right\} }{1-\max\left\{ \lambda_{2}(\bm{Q}^{\star}),-\lambda_{n}(\bm{Q}^{\star})\right\} }\geq\frac{\alpha}{\beta}\geq\frac{1}{2\kappa^{2}}.
\]
%
This together with the aforementioned  eigengap for $\bm{Q}^{\star}$
establishes the advertised result.



\paragraph{Proof of Lemma \ref{lemma:ranking-noise}.}

Let $\bm{D} \coloneqq \mathsf{diag}(\sqrt{\pi_1^{\star}},\cdots,\sqrt{\pi_n^{\star}})$. We have seen from the proof in Section~\ref{subsec:proof-DK-asym} (cf.~\eqref{eq:pi-norm-matrix-equation}) that
\[
	\|\bm{E}\|_{\bm{\pi}^{\star}} =\|\bm{D}\bm{E}\bm{D}^{-1}\| \leq  \|\bm{D}\|\, \|\bm{E}\| \, \|\bm{D}^{-1}\| = \frac{  \sqrt{\max_i \pi_i^{\star}} } {  \sqrt{\min_i \pi_i^{\star}} }  \|\bm{E}\|
	=  \sqrt{\kappa} \, \|\bm{E}\|,
\]
%
where $\kappa$ is defined in (\ref{eq:ranking-well-condition}).
Therefore, it suffices to bound $\|\bm{E}\|$.

% We first make the observation that
% %
% \begin{align*}
% \|\bm{E}\|_{\bm{\pi}^{\star}} & =\sup_{\bm{x}\neq\bm{0}}\frac{\|\bm{E}\bm{x}\|_{\bm{\pi}^{\star}}}{\|\bm{x}\|_{\bm{\pi}^{\star}}}=\sup_{\bm{x}\neq\bm{0}}\frac{\|\bm{D}\bm{E}\bm{D}^{-1}\bm{D}\bm{x}\|_{2}}{\|\bm{D}\bm{x}\|_{2}}=\sup_{\bm{v}\neq\bm{0}}\frac{\|\bm{D}\bm{E}\bm{D}^{-1}\bm{v}\|_{2}}{\|\bm{v}\|_{2}}\\
%  & =\|\bm{D}\bm{E}\bm{D}^{-1}\|,
% \end{align*}
%
% where $\bm{D} \coloneqq \mathsf{diag}(\sqrt{\pi_1^{\star}},\cdots,\sqrt{\pi_n^{\star}})$. Here, the second relation makes use of the identity $\|\bm{x}\|_{\bm{\pi}^{\star}} = \|\bm{D}\bm{x}\|_2$ (see the definition of the vector norm $\|\cdot\|_{\bm{\pi}^{\star}}$), the third relation replaces $\bm{D}\bm{x}$ with $\bm{v}$,  while the last one holds true due to the definition of the spectral norm.  As a result,
% %



By construction of $\bm{P}$ and $\bm{P}^{\star}$ (see \eqref{defn:ranking-P-empirical} and \eqref{defn:ranking-P-star}), we see that
%
\begin{equation}
	E_{i,j}=P_{i,j}-P_{i,j}^{\star}=\frac{1}{n} \big( y_{i,j}- \mathbb{E} [ y_{i,j}] \big) \label{eq:ranking-noise-off-diag}
\end{equation}
%
for any $i\neq j$.
In addition, for all $1\leq i\leq n$, it follows that
%
\begin{align}
E_{i,i} & =P_{i,i}-P_{i,i}^{\star}
%=\Big(1-\sum_{j:j\neq i}P_{i,j}\Big)-\Big(1-\sum_{j:j\neq i}P_{i,j}^{\star}\Big)
= -\sum_{j:j\neq i}E_{i,j}
  =-\frac{1}{n}\sum_{j:j\neq i}\big( y_{i,j}- \mathbb{E} [ y_{i,j}] \big) .\label{eq:ranking-noise-diag}
\end{align}
%
In view of these identities, we shall decompose the matrix $\bm{E}$
into three parts: the upper triangular part (denoted by $\bm{E}_{\mathsf{upper}}$),
the diagonal part (denoted by $\bm{E}_{\mathsf{diag}}$), and the lower triangular
part (denoted by $\bm{E}_{\mathsf{lower}}$). Clearly,  the triangle inequality gives
%
\begin{equation}
\|\bm{E}\|\leq\|\bm{E}_{\mathsf{upper}}\|+\|\bm{E}_{\mathsf{diag}}\|+\|\bm{E}_{\mathsf{lower}}\|.\label{eq:ranking-triangle}
\end{equation}
%
In the sequel, we deal with these three terms separately.

Let us start with the diagonal part $\bm{E}_{\mathsf{diag}}$.
In view of the definition of the spectral norm, we know that
%
\[
\|\bm{E}_{\mathsf{diag}}\|=\max_{1\leq i\leq n}|E_{i,i}|=\max_{1\leq i\leq n} \Big| \sum_{j:j\neq i}E_{i,j} \Big|,
\]
%
where the last relation arises from (\ref{eq:ranking-noise-diag}). Fix any
$i$, then it is easily seen that $\sum_{j:j\neq i}E_{i,j}$ is a sum of
independent zero-mean random variables $\{E_{i,j}\}$, which can be controlled via the Bernstein inequality.
Specifically, observe that
%
\[
\max_{j:j\neq i}|E_{i,j}|=\max_{j:j\neq i} \frac{1}{n} \big| y_{i,j}-\mathbb{E} [ y_{i,j}]   \big|
\leq\frac{1}{n} \eqqcolon B_1
\]
%
and,  in addition,
%
\[
	v_1\coloneqq\sum_{j:j\neq i}\mathbb{E}[E_{i,j}^{2}]=\frac{1}{n^{2}}\sum_{j:j\neq i}\mathsf{Var}\left(y_{i,j}\right)\leq\frac{1}{n},
\]
%
where the last inequality follows since the variance of a Bernoulli random variable is no larger than $1$.
Apply the Bernstein inequality (cf.~Corollary~\ref{thm:matrix-Bernstein-friendly}) and the union
bound over $1\leq i\leq n$ to demonstrate that
%
\begin{align*}
	\|\bm{E}_{\mathsf{diag}}\| & =\max_{1\leq i\leq n} \Big|\sum_{j:j\neq i}E_{i,j} \Big|
	\lesssim \sqrt{v_1\log n} + B_1 \log n   \\
	& \lesssim  \sqrt{\frac{\log n}{n}}+\frac{\log n}{n}\asymp\sqrt{\frac{\log n}{n}}
\end{align*}
%
holds with probability at least $1-O(n^{-8})$.
% This finishes the upper bound on the diagonal part.

We now move on to the upper triangular part $\bm{E}_{\mathsf{upper}}$,
whose entries $\{E_{i,j}\}_{i<j}$ are independent. Invoking Theorem~\ref{thm:tighter-spectral-normal-ramon} (see the remark about asymmetric version right after Theorem~\ref{thm:tighter-spectral-normal-ramon}) with the bounds on $B_1$ and $v_1$ established above,
we arrive at
%
\[
	\|\bm{E}_{\mathsf{upper}}\|\lesssim \sqrt{v_1 } + B_1 \log n \lesssim \sqrt{\frac{1}{n}} + \frac{\log n}{n} \asymp  \sqrt{\frac{1}{n}}
\]
with probability at least $1-O(n^{-8})$. Similar arguments lead to
the same upper bound on $\|\bm{E}_{\mathsf{lower}}\|$, which
we omit  for brevity.

Substituting the upper bounds on $\|\bm{E}_{\mathsf{diag}}\|$, $\|\bm{E}_{\mathsf{upper}}\|$ and $\|\bm{E}_{\mathsf{lower}}\|$ into (\ref{eq:ranking-triangle}), we immediately establish
the desired bound.

%This entrywise independence renders Theorem~\ref{thm:tighter-spectral-normal-ramon} particularly useful, with the only caveat that Theorem~\ref{thm:tighter-spectral-normal-ramon} applies to symmetric matrices.
%To address this issue, we introduce the symmetric dilation $\mathcal{S}(\bm{E}_{\mathsf{upper}})$ of $\bm{E}_{\mathsf{upper}}$:
%\[
%\mathcal{S}(\bm{E}_{\mathsf{upper}})\coloneqq\left[\begin{array}{cc}
%\bm{0} & \bm{E}_{\mathsf{upper}}\\
%\bm{E}_{\mathsf{upper}}^{\top} & \bm{0}
%\end{array}\right],
%\]
%which has the desired symmetry. But more importantly, $\|\mathcal{S}(\bm{E}_{\mathsf{upper}})\| = \|\bm{E}_{\mathsf{upper}}\|$.

%In light of (\ref{eq:ranking-noise-off-diag}), it can be written as
%%
%\[
%\bm{E}_{\mathsf{upper}}=\sum_{i<j}E_{i,j}\bm{e}_{i}\bm{e}_{j}^{\top}=  \sum_{i<j} \frac{1}{n} \big(y_{i,j}-\mathbb{E} [ y_{i,j}] \big)\bm{e}_{i}\bm{e}_{j}^{\top}
%\eqqcolon \sum_{i<j} \bm{Z}_{i,j}
%\]
%%
%which is composed of independent zero-mean random matrices $\{\bm{Z}_{i,j}\}$. These matrices obey
%%
%\[
%	\max_{i<j}\|\bm{Z}_{i,j}\|\leq\max_{i<j}|E_{i,j}|\leq \frac{1}{n} \eqqcolon B_{2}
%\]
%%
%and, in addition,
%%
%\begin{align*}
%v_{2} & \coloneqq\max\Bigg\{ \Big\|\sum_{i<j}\mathbb{E}\big[\bm{Z}_{i,j}\bm{Z}_{i,j}^{\top}\big]\Big\|,\,\Big\|\sum_{i<j}\mathbb{E}\big[\bm{Z}_{i,j}^{\top}\bm{Z}_{i,j}\big]\Big\|\Bigg\} \\
% & =\max\Bigg\{ \Big\|\sum_{i<j}\mathbb{E}\big[E_{i,j}^{2}\bm{e}_{i}\bm{e}_{i}^{\top}\big]\Big\|,\,\Big\|\sum_{i<j}\mathbb{E}\big[E_{i,j}^{2}\bm{e}_{j}\bm{e}_{j}^{\top}\big]\Big\|\Bigg\} \\
%	& \leq \frac{1}{n^2} \max\Bigg\{ \Big\|\sum_{i<j}\bm{e}_{i}\bm{e}_{i}^{\top}\Big\|,\,\Big\|\sum_{i<j}\bm{e}_{j}\bm{e}_{j}^{\top}\Big\|\Bigg\}
%	\leq \frac{1}{n^2}\|n\bm{I}\|= \frac{1}{n}.
%\end{align*}
%%
%Here, the last line again invokes the condition $\mathsf{Var}(y_{i,j})\leq 1$.
